\documentclass[12pt,a4paper]{article}
\usepackage{amsmath, amssymb, amsthm}
\usepackage{geometry}
\usepackage{graphicx}
\usepackage{hyperref}
\usepackage{xcolor}
\usepackage{booktabs}
\usepackage{algorithm}
\usepackage{algpseudocode}
\usepackage{tcolorbox}
\usepackage{enumitem}
\usepackage{mathtools}
\usepackage{tikz}
\usetikzlibrary{graphs, positioning}

\geometry{margin=1in}

\hypersetup{
    colorlinks=true,
    linkcolor=blue!70!black,
    citecolor=green!60!black,
    urlcolor=cyan!80!black
}

\tcbuselibrary{skins, breakable}
\newtcolorbox{failbox}[1][]{colback=red!5!white, colframe=red!60!black,
    fonttitle=\bfseries, title=Failure Case, #1}
\newtcolorbox{improvebox}[1][]{colback=blue!5!white, colframe=blue!60!black,
    fonttitle=\bfseries, title=Proposed Improvement, #1}
\newtcolorbox{crazybox}[1][]{colback=purple!5!white, colframe=purple!60!black,
    fonttitle=\bfseries, title=Crazy Idea, #1}

\title{
    \textbf{Failure Cases and Radical Improvements for} \\
    \textbf{Optimization-Based Graph Isomorphism Detection} \\[0.5em]
    \large A Critical Analysis of the Continuous Relaxation Approach
}
\author{Graph Isomorphism Research}
\date{\today}

\begin{document}

\maketitle
\tableofcontents
\newpage

%% ============================================================
\section{Introduction}
%% ============================================================

The continuous relaxation approach to graph isomorphism (GI) is an audacious idea: replace
the combinatorial search over all $n!$ permutations with a smooth optimization over the
Birkhoff polytope $\mathcal{B}_n$ of doubly stochastic matrices. The current objective is:

\begin{equation}
    Z^* = \min_{X \in \mathcal{B}n} \underbrace{\langle C, X \rangle}{\text{degree compat}}
    + \underbrace{\|AX - XB\|F^2}{Q(X)}
    + \lambda_T \underbrace{\|A^2 X - X B^2\|F^2}{T(X)}
    + \lambda_S \underbrace{\|U_A - X U_B\|F^2}{S(X)}
    \label{eq:main}
\end{equation}

with isomorphism index $I = e^{-\lambda Z^*}$, where $I \approx 1$ signals isomorphism and $I \ll 1$
signals non-isomorphism.

\medskip
While elegant, this formulation has \textbf{provably exploitable failure modes}. This report
systematically enumerates them, provides concrete graph examples, and proposes—sometimes
radical—improvements to push toward a complete and correct solver.

%% ============================================================
\section{Failure Case 1: The Threshold Problem}
%% ============================================================

\subsection{Description}

The decision rule $I \approx 1 \Rightarrow \text{isomorphic}$ is inherently fuzzy.
The mapping $I = e^{-\lambda Z^*}$ with default $\lambda = 0.1$ is a \textbf{heuristic with
no theoretical threshold guarantee}.

\begin{failbox}[title=Failure Case 1 — Ambiguous Threshold]
For isomorphic graphs, the QP relaxation is \emph{not always tight}: $X^*$ may be
fractional even when a true permutation exists. This causes $Z^* > 0$ for genuinely
isomorphic graphs, pushing $I$ below 1 and potentially misclassifying them as
\textbf{non-isomorphic}.

\medskip
\textbf{Concrete Example:} Any graph with non-trivial automorphism group (e.g., the complete
graph $K_n$, cycle $C_n$, or Petersen graph). The solver finds infinitely many optimal
$X^*$ (all convex combinations of valid permutations), none of which is binary, so
$Z^* > 0$ even though the graphs are isomorphic.
\end{failbox}

\subsection{Why It Fails}

The Birkhoff--von Neumann theorem guarantees that the vertices of $\mathcal{B}_n$ are
permutation matrices. However, when $Q(X) = 0$ at multiple permutations simultaneously
(i.e., the graph has automorphisms), the QP objective has a \emph{flat region} over a
face of the polytope. Interior-point solvers converge to the analytic center of this face,
which is fractional.

\subsection{Proposed Fix}

\begin{improvebox}[title=Hungarian Rounding + Exact Residual Verification]
After solving the QP to get $X^*$:
\begin{enumerate}
    \item Apply the Hungarian algorithm to extract the nearest permutation:
    \[
        P^* = \arg\min_{P \in \Pi_n} \|P - X^*\|_F
    \]
    \item Compute the exact structural residual:
    \[
        r = \|AP^* - P^*B\|_F
    \]
    \item Decide: $r < \varepsilon \Rightarrow$ \textbf{isomorphic}, else \textbf{not isomorphic}.
\end{enumerate}
This converts the fuzzy score into a \textbf{binary, provably correct} answer whenever the
relaxation is tight.
\end{improvebox}

%% ============================================================
\section{Failure Case 2: Strongly Regular Graphs}
%% ============================================================

\subsection{Description}

A strongly regular graph $\text{srg}(n, k, \lambda, \mu)$ has every vertex with degree $k$,
every adjacent pair with $\lambda$ common neighbors, and every non-adjacent pair with
$\mu$ common neighbors.

\begin{failbox}[title=Failure Case 2 — Strongly Regular Graphs]
Two non-isomorphic strongly regular graphs with the \emph{same parameters}
$(n, k, \lambda, \mu)$ will have:
\begin{itemize}
    \item Identical degree sequences $\Rightarrow$ $C_{ij}$ is \textbf{uniform} (no signal)
    \item $\text{tr}(A^2) = \text{tr}(B^2) = nk$ $\Rightarrow$ 1-hop commutator provides no discrimination
    \item $\text{tr}(A^3) = \text{tr}(B^3) = nk\lambda$ $\Rightarrow$ 2-hop commutator also provides no discrimination
    \item Identical Laplacian spectrum $\Rightarrow$ $S(X)$ is minimized at $X = \frac{1}{n}\mathbf{1}\mathbf{1}^\top$ (the center of $\mathcal{B}_n$)
\end{itemize}
\textbf{Result}: All four terms of $Z^*$ are simultaneously fooled. The objective
\emph{cannot distinguish} the two graphs.

\medskip
\textbf{Concrete Example:} The two non-isomorphic $\text{srg}(16, 6, 2, 2)$ graphs (the
$4 \times 4$ rook's graph vs.\ the Shrikhande graph). All invariants in the current
objective are identical.
\end{failbox}

\subsection{Why It Fails}

The objective $Z^*$ only encodes information up to \emph{2-hop walk counts}. Strongly regular
graphs are defined precisely to make all such local statistics identical. The spectral term
also fails because cospectral mates exist for all srg families.

\subsection{Proposed Fix}

\begin{improvebox}[title=Higher-Order Commutator Cascade]
Add commutator terms for $k = 3, 4, \ldots, K$:
\[
    Z^* \mathrel{+}= \sum_{k=3}^{K} \lambda_k \|A^k X - X B^k\|_F^2
\]
For the Shrikhande vs.\ rook's graph example, $k=4$ suffices to distinguish them.
More generally, set $K = \lceil \text{diam}(G) \rceil$.
\end{improvebox}

\begin{crazybox}[title=Crazy Idea — Graph Zeta Function Alignment]
Use the Ihara zeta function $\zeta_G(u)^{-1} = \det(I - Au + (D - I)u^2)$ as an
additional invariant. Two graphs are zeta-equivalent only if their spectra of the
\emph{edge adjacency operator} match — strictly stronger than Laplacian cospectrality.
Add the penalty:
\[
    \Phi(X) = \left\| \log \zeta_{G_1}(u) - \log \zeta_{G_2}(u) \right\|_{L^2}
\]
evaluated at a grid of $u$ values. This is cheap to compute and extraordinarily
discriminative for regular graphs.
\end{crazybox}

%% ============================================================
\section{Failure Case 3: Eigenspace Degeneracy in the Spectral Term}
%% ============================================================

\subsection{Description}

The spectral alignment term is:
\[
    S(X) = \|U_A - X U_B\|_F^2
\]
where $U_A, U_B$ are matrices of Laplacian eigenvectors ordered by eigenvalue.

\begin{failbox}[title=Failure Case 3 — Degenerate Eigenspaces]
When a Laplacian has a repeated eigenvalue of multiplicity $m > 1$, the corresponding
$m$-dimensional eigenspace has \textbf{no canonical basis}. Any orthonormal basis
for this subspace is equally valid. The sign-normalization heuristic in the current
implementation \emph{does not resolve this ambiguity for $m > 1$}.

\medskip
\textbf{Concrete Example:} All $k$-regular graphs have Laplacian eigenvalue 0 with
multiplicity equal to the number of connected components, and eigenvalue $k$ with
multiplicity equal to the bipartite structure. For the complete graph $K_n$, the
eigenvalue $n$ has multiplicity $n-1$, making $U_A$ and $U_B$ essentially random
rotations of each other.

\medskip
\textbf{Result}: $S(X)$ provides \textbf{misleading gradient information}, pushing
the solver away from correct solutions.
\end{failbox}

\subsection{Proposed Fix}

\begin{improvebox}[title=Eigenvalue-Weighted Grassmannian Alignment]
Instead of comparing specific eigenvectors, compare the \emph{eigenspaces} as points
on the Grassmannian manifold. For each distinct eigenvalue cluster, compute the projection
operator:
\[
    \Pi_A^{(\ell)} = U_A^{(\ell)} (U_A^{(\ell)})^\top, \quad
    \Pi_B^{(\ell)} = U_B^{(\ell)} (U_B^{(\ell)})^\top
\]
and penalize:
\[
    S'(X) = \sum_\ell w_\ell \left\| \Pi_A^{(\ell)} - X \Pi_B^{(\ell)} X^\top \right\|_F^2
\]
where $w_\ell = |\lambda_\ell|$ weights by eigenvalue magnitude. This is \textbf{rotation-invariant
within eigenspaces} and correctly handles all degeneracies.
\end{improvebox}

%% ============================================================
\section{Failure Case 4: The Unimplemented Entropy Regularizer}
%% ============================================================

\subsection{Description}

The paper proposes the entropy continuation term:
\[
    H(X) = -\varepsilon \sum_{i,j} X_{ij} \log X_{ij}
\]
but states it is \textbf{not implemented} because it breaks the QP structure.

\begin{failbox}[title=Failure Case 4 — Cold-Start Local Minima]
Without entropy regularization, the QP is initialized at a point in $\mathcal{B}_n$
and may converge to a \textbf{local minimum} of the non-convex landscape (note:
$Q(X)$, $T(X)$, $S(X)$ are all non-convex in $X$, even though the feasible set
is convex). The solver finds the nearest stationary point, which may not be the global
optimum.

\medskip
\textbf{Concrete Example:} Two isomorphic graphs where the initialization is close to
a \emph{wrong} permutation (e.g., an automorphism). The solver gets trapped near this
wrong permutation, returning $Z^* > 0$ when the true minimum is $Z^* = 0$.
\end{failbox}

\subsection{Proposed Fix}

\begin{improvebox}[title=Sinkhorn Warm-Start + Graduated Non-Convexity]
Use a two-phase approach:
\begin{enumerate}
    \item \textbf{Phase 1 — Sinkhorn Entropic Relaxation:}
    Solve the regularized assignment problem:
    \[
        X^{(0)} = \text{Sinkhorn}\!\left(e^{-C/\varepsilon}, \text{max\_iter}=200\right)
    \]
    with temperature $\varepsilon$ starting large (diffuse) and annealing to 0.
    \item \textbf{Phase 2 — QP Warm Start:} Initialize Gurobi with $X^{(0)}$ instead of
    a uniform or zero start.
    \item \textbf{Phase 3 — Multi-restart:} Run $k$ random perturbations of $X^{(0)}$,
    solve each QP, take the minimum $Z^*$.
\end{enumerate}
\end{improvebox}

\begin{crazybox}[title=Crazy Idea — Quantum-Inspired Annealing via QUBO]
Reformulate as a Quadratic Unconstrained Binary Optimization (QUBO):
\[
    \min_{P \in \{0,1\}^{n^2}} P^\top M P + \text{constraint penalties}
\]
and solve with \textbf{simulated quantum annealing} (D-Wave or classically via path-integral
Monte Carlo). The tunneling effect escapes local minima that trap classical gradient methods.
For $n \leq 50$, this is feasible on current quantum annealers.
\end{crazybox}

%% ============================================================
\section{Failure Case 5: Weak Node Compatibility Matrix $C$}
%% ============================================================

\subsection{Description}

The current cost matrix uses only degree and 1-hop degree profile:
\[
    C_{ij} = \alpha |d^A_i - d^B_j| + \beta |s^A_i - s^B_j|
\]

\begin{failbox}[title=Failure Case 5 — Uniform $C$ Matrix]
For $k$-regular graphs, all nodes have the same degree $k$ and the same degree-profile
sum $k^2$, making $C_{ij} = 0$ for all $(i,j)$.

\medskip
\textbf{Result}: The linear term $\langle C, X \rangle = 0$ everywhere — it contributes
\textbf{zero gradient information}. The entire computational burden falls on the quadratic
terms $Q, T, S$, which (as shown in Failure Cases 2 and 3) also fail on regular graphs.

\medskip
\textbf{Concrete Example:} The Petersen graph (3-regular, 10 nodes). Every node has
degree 3, $s_i = 9$, so $C = \mathbf{0}$ and the linear term vanishes completely.
\end{failbox}

\subsection{Proposed Fix}

\begin{improvebox}[title=WL-Enriched Node Compatibility via $k$-hop Fingerprints]
Replace $C$ with a multi-scale fingerprint:
\[
    C_{ij} = \sum_{k=1}^{K} \gamma_k \left\| \phi^{(k)}_i(A) - \phi^{(k)}_j(B) \right\|_2
\]
where $\phi^{(k)}_i(G)$ is the \emph{sorted $k$-hop degree sequence} of node $i$ in graph $G$:
\[
    \phi^{(k)}i(G) = \text{sort}\!\left\{ \deg{G^k}(u) : u \in \mathcal{N}^k_i \right\}
\]
For $K = 4$, this distinguishes nodes in all known problematic graph families at
negligible computational cost.
\end{improvebox}

%% ============================================================
\section{Failure Case 6: Numerical Precision for Large $n$}
%% ============================================================

\begin{failbox}[title=Failure Case 6 — Floating Point Catastrophic Cancellation]
The Frobenius norm $\|AX - XB\|_F^2$ involves a difference of two dense matrix products.
For large $n$ (even $n \approx 50$) and graphs with entries near 0 and 1, catastrophic
cancellation causes the computed residual to have relative error up to $O(n \cdot \epsilon_{\text{mach}})$.

\medskip
\textbf{Result}: The solver declares $Z^* \approx 0$ for \textbf{non-isomorphic} graphs
due to floating-point error, producing false positives.
\end{failbox}

\begin{improvebox}[title=Compensated Summation + Integer Certificate]
\begin{enumerate}
    \item Use Kahan compensated summation when computing $\|R\|_F^2$.
    \item After rounding to $P^$ via Hungarian, verify $AP^ = P^*B$ \textbf{in integer arithmetic}
    (adjacency matrices are 0/1). This gives a \textbf{provably exact} binary answer.
\end{enumerate}
\end{improvebox}

%% ============================================================
\section{Crazy Improvements for Near-Complete Correctness}
%% ============================================================

\subsection{Crazy Idea 1: Homomorphism Density Tensor}

\begin{crazybox}[title=Homomorphism Density Matching]
The \emph{homomorphism density} $t(F, G)$ counts the fraction of maps from a small
``pattern'' graph $F$ into $G$ that are graph homomorphisms. Two graphs are isomorphic
only if $t(F, G_1) = t(F, G_2)$ for \emph{all} $F$ (Lov\'asz's theorem). Add:
\[
    \mathcal{H}(X) = \sum_{F \in \mathcal{F}} w_F \left| t(F, G_1) - t(F, G_2) \right|^2
\]
For $\mathcal{F} = \{K_2, K_3, K_4, C_4, P_3\}$, this is computable in polynomial
time and provides a \emph{theoretically grounded} completeness guarantee as $|\mathcal{F}| \to \infty$.
\end{crazybox}

\subsection{Crazy Idea 2: Probabilistic Hash Certificate}

\begin{crazybox}[title=Random Polynomial Identity Testing]
After finding $P^*$, use the Schwartz--Zippel lemma: evaluate the polynomial
\[
    f(t) = \det(A - t \cdot I) - \det(P^{\top} B P^ - t \cdot I)
\]
at $K$ random points $t_1, \ldots, t_K \in \mathbb{F}_p$ (a large prime field). If
$f(t_k) = 0$ for all $k$, declare isomorphic with error probability $\leq n/p$ per test.
For $p > n \cdot 2^{128}$ and $K = 10$, this gives essentially zero false positive rate
with $O(n^3)$ cost.
\end{crazybox}

\subsection{Crazy Idea 3: Neural Residual Oracle}

\begin{crazybox}[title=Neural Residual Oracle for Hard Cases]
Train a GNN on the \emph{hard cases} (strongly regular graphs, cospectral pairs) to
serve as a residual oracle. The pipeline becomes:
\[
    \text{Answer} =
    \begin{cases}
        \text{isomorphic} & \text{if } Z^* < \varepsilon_1 \text{ and } \|AP^* - P^*B\|_F = 0 \\
        \text{not isomorphic} & \text{if } Z^* > \varepsilon_2 \\
        \text{GNN Oracle}(G_1, G_2) & \text{if } \varepsilon_1 \leq Z^* \leq \varepsilon_2
    \end{cases}
\]
The oracle is only invoked for the ambiguous middle zone, keeping average cost low
while handling hard cases gracefully.
\end{crazybox}

%% ============================================================
\section{Complete Improved Objective}
%% ============================================================

Combining all non-crazy improvements, the proposed enhanced objective is:

\begin{equation}
    \boxed{
    Z^* = \min_{X \in \mathcal{B}_n}
    \underbrace{\langle C^{\text{WL}}, X \rangle}_{\text{WL node compat}}
    + \sum_{k=1}^{K} \lambda_k \underbrace{\|A^k X - X B^k\|F^2}{k\text{-hop commutator}}
    + \lambda_S \underbrace{S'(X)}_{\text{Grassmannian spectral}}
    + \lambda_\Phi \underbrace{\Phi(X)}_{\text{Zeta function}}
    }
    \label{eq:improved}
\end{equation}

with the post-processing pipeline:
\begin{enumerate}
    \item Sinkhorn warm start (Phase 1)
    \item Solve \eqref{eq:improved} with Gurobi (Phase 2)
    \item Hungarian rounding to $P^*$ (Phase 3)
    \item Integer arithmetic verification $\|AP^* - P^*B\|_F \stackrel{?}{=} 0$ (Phase 4)
    \item If ambiguous: invoke Neural Residual Oracle (Phase 5)
\end{enumerate}

%% ============================================================
\section{Summary of Failure Cases and Fixes}
%% ============================================================

\begin{table}[h!]
\centering
\renewcommand{\arraystretch}{1.4}
\begin{tabular}{@{}p{3.5cm} p{4cm} p{4cm} c@{}}
\toprule
\textbf{Failure Case} & \textbf{Graph Family} & \textbf{Fix} & \textbf{Cost} \\
\midrule
Fuzzy threshold / fractional $X^*$ & All graphs with automorphisms ($K_n$, $C_n$, Petersen) & Hungarian rounding + integer verify & $O(n^3)$ \\
\midrule
All 4 terms fooled simultaneously & Strongly regular graphs (Shrikhande vs.\ rook's) & $k$-hop commutators $k \leq 6$ & $O(Kn^3)$ \\
\midrule
Spectral term misleading & $k$-regular graphs, $K_n$ & Grassmannian eigenspace alignment & $O(n^3)$ \\
\midrule
Cold-start local minima & Any graph, wrong initialization & Sinkhorn warm-start + multi-restart & $O(n^2 \log n)$ \\
\midrule
Zero gradient from $C$ & Regular graphs (Petersen, $C_n$) & WL $k$-hop fingerprint & $O(Kn^2)$ \\
\midrule
Floating-point false positives & Large $n \gtrsim 50$ & Integer arithmetic certificate & $O(n^2)$ \\
\bottomrule
\end{tabular}
\caption{Summary of failure cases, affected graph families, proposed fixes, and their computational cost.}
\end{table}

%% ============================================================
\section{Conclusion}
%% ============================================================

The continuous relaxation approach to graph isomorphism is a genuinely bold idea. Its
current formulation is powerful for generic graphs but has \emph{provable} blind spots
for structured families — particularly strongly regular graphs, regular graphs, and
graphs with large automorphism groups. These are not edge cases; they are the
\emph{hardest} instances in competitive GI solvers.

The path to near-completeness requires three layers:
\begin{enumerate}
    \item \textbf{Richer invariants}: Move from 2-hop to $K$-hop commutators, WL fingerprints,
    and zeta function alignment.
    \item \textbf{Better optimization}: Sinkhorn warm starts, graduated non-convexity,
    and multi-restart strategies.
    \item \textbf{Exact post-processing}: Replace the fuzzy isomorphism index with
    Hungarian rounding followed by provably correct integer arithmetic verification.
\end{enumerate}

With these modifications, the optimization-based pipeline becomes not just a heuristic
but a \emph{practically complete} algorithm — one that fails only on graph instances
that defeat every known polynomial-time method.

\end{document}